\documentclass[../../main.tex]{subfiles}

\graphicspath{{../../image/}{../../figure/}}

\begin{document}

\section{习题}\label{section:ch20-习题}

\begin{exercise}\label{exercise:ch20-01}
（★）使用\eqref{align:ch20-递推}，以 $\bm{z}_{t-1}$ 的均值和协方差为变量，写出 $\bm{z}_t$ 的均值和协方差。证明当 $0<\beta_t<1$ 时，$\bm{z}_t$ 的均值比 $\bm{z}_{t-1}$ 的均值更接近于零，并且 $\bm{z}_t$ 的协方差比 $\bm{z}_{t-1}$ 的协方差更接近于单位矩阵 $\bm{I}$。
\end{exercise}

\begin{exercise}\label{exercise:ch20-02}
（★）证明\eqref{align:ch20-第一步}可以等价地表示为\eqref{align:ch20-一步分布}。
\end{exercise}

\begin{exercise}\label{exercise:ch20-03}
（★★★）使用归纳法证明扩散模型前向过程的 $\bm{z}_t$ 的边缘分布由\eqref{align:ch20-扩散核}给出，其中 $\alpha_t$ 由\eqref{align:ch20-alpha}定义。首先验证当 $t=1$ 时\eqref{align:ch20-扩散核}成立。然后假设对于某个特定的 $t$ 值，\eqref{align:ch20-扩散核}也成立，并且推导出 $t+1$ 时的相应结果。为此，最简单的方法是使用\eqref{align:ch20-递推}来写前向过程，并利用“两个独立 Gauss 随机变量之和仍是 Gauss 分布、且均值与协方差可加”这一结果。
\end{exercise}

\begin{exercise}\label{exercise:ch20-04}
（★）利用\eqref{align:ch20-扩散核}，其中 $\alpha_t$ 由\eqref{align:ch20-alpha}定义，证明在极限 $T\to\infty$ 下可以得到\eqref{align:ch20-极限}。
\end{exercise}

\begin{exercise}\label{exercise:ch20-05}
（★★）考虑两个独立的随机变量 $\bm{a}$ 和 $\bm{b}$ 以及一个固定的标量 $\lambda$，证明
\begin{align}
\operatorname{Cov}[\bm{a}+\bm{b}]&=\operatorname{Cov}[\bm{a}]+\operatorname{Cov}[\bm{b}],
\label{align:ch20-协方差可加}\\
\operatorname{Cov}[\lambda\bm{a}]&=\lambda^2\operatorname{Cov}[\bm{a}].
\label{align:ch20-协方差缩放}
\end{align}
然后进一步证明：如果 $\bm{z}_{t-1}$ 的分布具有零均值和单位协方差，那么无论 $\beta_t$ 的值如何，\eqref{align:ch20-递推}定义的 $\bm{z}_t$ 的分布都具有零均值和单位协方差。
\end{exercise}

\begin{exercise}\label{exercise:ch20-06}
（★★★）使用“配方法”这一技术从 Bayes 定理\eqref{align:ch20-条件逆}中推导出\eqref{align:ch20-逆条件}。首先注意\eqref{align:ch20-条件逆}右边分子中的两项，它们分别由\eqref{align:ch20-转移核}和\eqref{align:ch20-扩散核}给出，并且它们都采用关于 $\bm{z}_{t-1}$ 的二次函数的指数形式。因此所需的分布是 Gauss 分布，我们只需要求出其均值和协方差即可。为此，请只考虑指数中依赖于 $\bm{z}_{t-1}$ 的项，注意两个指数的乘积等于这两个指数之和的指数形式。把所有关于 $\bm{z}_{t-1}$ 的二次项以及关于 $\bm{z}_{t-1}$ 的一次项收集起来，然后重新整理成 $(\bm{z}_{t-1}-\bm{m}_t)^{\mathsf T}\bm{S}_t^{-1}(\bm{z}_{t-1}-\bm{m}_t)$ 的形式。最后通过观察，求出 $\bm{m}_t(\bm{x},\bm{z}_t)$ 和 $\bm{S}_t$ 的表达式。注意，与 $\bm{z}_{t-1}$ 无关的加法项可以忽略不计。
\end{exercise}

\begin{exercise}\label{exercise:ch20-07}
（★★★）证明扩散模型中正向加噪过程的条件分布 $q(\bm{z}_t\mid\bm{z}_{t-1})$ 的逆，在噪声方差较小时可以近似为 Gauss 分布。考虑由 Bayes 定理\eqref{align:ch20-Bayes逆}给出的逆条件分布 $q(\bm{z}_{t-1}\mid\bm{z}_t)$，而正向分布 $q(\bm{z}_t\mid\bm{z}_{t-1})$ 由\eqref{align:ch20-转移核}给出。对\eqref{align:ch20-Bayes逆}的两边取对数，然后以 $\bm{z}_t$ 为中心对 $q(\bm{z}_{t-1})$ 进行 Taylor 展开，证明对于小的噪声方差 $\beta_t$，分布 $q(\bm{z}_{t-1}\mid\bm{z}_t)$ 可以近似为均值为 $\bm{z}_t$、协方差为 $\beta_t\bm{I}$ 的 Gauss 分布。求出关于均值和协方差的最低阶修正项的表达式，将其表示为 $\beta_t$ 的幂次展开形式。
\end{exercise}

\begin{exercise}\label{exercise:ch20-08}
（★★）通过将概率的乘积法则\eqref{align:ch20-乘积法则}代入扩散模型的 ELBO 定义\eqref{align:ch20-ELBO定义}中，并利用 KL 散度的定义\eqref{align:ch20-KL定义}，验证对数似然函数可以写成下界和 KL 散度之和的形式\eqref{align:ch20-ELBO恒等}。
\end{exercise}

\begin{exercise}\label{exercise:ch20-09}
（★★）验证扩散模型的 ELBO［由\eqref{align:ch20-中间下界}给出］可以写成\eqref{align:ch20-KL分解}的形式，其中 KL 散度由\eqref{align:ch20-KL定义}定义。
\end{exercise}

\begin{exercise}\label{exercise:ch20-10}
（★★）在推导\eqref{align:ch20-KL分解}给出的扩散模型的 ELBO 时，我们省略了\eqref{align:ch20-展开下界}右边的第一项和第三项，因为它们与 $\bm{w}$ 无关。同样，我们还省略了\eqref{align:ch20-对数比}右边的第二项，因为它也与 $\bm{w}$ 无关。证明如果保留所有这些省略的项，ELBO 中就会增加一个额外的项
\begin{align}
\operatorname{KL}\big(q(\bm{z}_T\mid\bm{x})\,\|\,p(\bm{z}_T)\big).
\label{align:ch20-额外KL}
\end{align}
请注意加噪过程是这样构造的：设法使分布 $q(\bm{z}_T\mid\bm{x})$ 等于 Gauss 分布 $\mathcal{N}(\bm{z}_T\mid\bm{0},\bm{I})$。同样，分布 $p(\bm{z}_T)$ 定义为等同于 $\mathcal{N}(\bm{z}_T\mid\bm{0},\bm{I})$，因此\eqref{align:ch20-额外KL}中的两个分布相同，KL 散度为 $0$。
\end{exercise}

\begin{exercise}\label{exercise:ch20-11}
（★★）利用\eqref{align:ch20-逆条件}所代表的分布 $q(\bm{z}_{t-1}\mid\bm{z}_t,\bm{x})$ 和\eqref{align:ch20-反向过程}所代表的分布 $p(\bm{z}_{t-1}\mid\bm{z}_t,\bm{w})$，证明出现在\eqref{align:ch20-KL分解}中的一致项中的 KL 散度由\eqref{align:ch20-一致项}给出。
\end{exercise}

\begin{exercise}\label{exercise:ch20-12}
（★★）将\eqref{align:ch20-x反解}代入\eqref{align:ch20-逆均值}中，重写均值 $\bm{m}_t(\bm{x},\bm{z}_t)$，使其能够以原始数据向量 $\bm{x}$ 和噪声 $\bm{\varepsilon}$ 的形式来表示，参见\eqref{align:ch20-均值重写}，其中 $\bm{\varepsilon}$ 由\eqref{align:ch20-直接采样}定义。
\end{exercise}

\begin{exercise}\label{exercise:ch20-13}
（★★）证明扩散模型的 ELBO 中的重建项［\eqref{align:ch20-重构项}］可以写成\eqref{align:ch20-重构噪声}的形式。考虑用\eqref{align:ch20-网络均值}替换 $\bm{\mu}(\bm{z}_1,\bm{w},1)$，并用\eqref{align:ch20-第一步}替换 $\bm{x}$，然后利用 $\alpha_1=(1-\beta_1)$，后者可以由\eqref{align:ch20-alpha}得出。
\end{exercise}

\begin{exercise}\label{exercise:ch20-14}
（★）得分函数是由 $\bm{s}(\bm{x})=\nabla_{\bm{x}}\ln p(\bm{x}\mid\bm{w})$ 定义的，因此它是一个与输入向量 $\bm{x}$ 维度相同的向量。考虑一个矩阵，其元素由下式给出：
\begin{align}
M_{ij}=\frac{\partial s_i}{\partial x_j}-\frac{\partial s_j}{\partial x_i}.
\label{align:ch20-反对称}
\end{align}
证明如果通过取梯度 $\bm{s}=\nabla_{\bm{x}}\phi(\bm{x})$ 来定义得分函数，其中 $\phi(\bm{x})$ 是具有单一输出变量的神经网络的输出，那么对于所有 $i,j$，矩阵元素 $M_{ij}$ 都为 $0$。请注意，如果得分函数 $\bm{s}(\bm{x})=\nabla_{\bm{x}}\ln p(\bm{x}\mid\bm{w})$ 直接用一个深度神经网络来表示，并且输入与输出的数量相同，那么只有对角线矩阵元素 $M_{ii}=0$，因此网络的输出不一定是某标量函数的梯度。
\end{exercise}

\begin{exercise}\label{exercise:ch20-15}
（★★）考虑使用一个深度神经网络 $\bm{s}(\bm{x},\bm{w})$ 来表示得分函数［由\eqref{align:ch20-得分}定义］，其中 $\bm{x}$ 和 $\bm{s}$ 的维度为 $D$。比较以下两种情况在计算得分时的计算复杂度：一种是具有 $D$ 个输出且直接用来表示得分函数的网络，另一种是计算单一标量函数 $\phi(\bm{x},\bm{w})$，并通过自动微分间接计算得分函数的网络。证明通常后一种方法的计算成本更高。
\end{exercise}

\begin{exercise}\label{exercise:ch20-16}
（★★★）我们无法直接最小化得分函数［\eqref{align:ch20-得分损失}］，因为我们不知道真实数据密度 $p(\bm{x})$ 的函数形式，因此无法写出得分函数 $\nabla_{\bm{x}}\ln p(\bm{x})$ 的表达式。然而，通过使用分部积分（Hyvärinen, 2005），我们可以将\eqref{align:ch20-得分损失}改写为
\begin{align}
J(\bm{w})=\int\Big\{\nabla\cdot\bm{s}(\bm{x},\bm{w})+\frac{1}{2}\big\|\bm{s}(\bm{x},\bm{w})\big\|^2\Big\}p(\bm{x})\,\mathrm{d}\bm{x}+\mathrm{const},
\label{align:ch20-得分损失改写}
\end{align}
其中常数项 $\mathrm{const}$ 不依赖于网络参数 $\bm{w}$，散度 $\nabla\cdot\bm{s}(\bm{x},\bm{w})$ 定义如下：
\begin{align}
\nabla\cdot\bm{s}=\sum_{i=1}^{D}\frac{\partial s_i}{\partial x_i}=\sum_{i=1}^{D}\frac{\partial^2\ln p(\bm{x})}{\partial x_i^2},
\label{align:ch20-散度}
\end{align}
其中 $D$ 是 $\bm{x}$ 的维度。首先展开\eqref{align:ch20-得分损失}中的平方项，注意涉及 $\|\bm{s}(\bm{x},\bm{w})\|^2$ 的项已经出现在\eqref{align:ch20-得分损失}中，而涉及 $\|\bm{s}_p\|^2$ 的项可以吸收到加法常数中，这里我们定义 $\bm{s}_p=\nabla\ln p(\bm{x})$。考虑下式对两个函数乘积的导数：
\begin{align}
\frac{\mathrm{d}}{\mathrm{d}x}\big\{p(x)g(x)\big\}=\frac{\mathrm{d}p(x)}{\mathrm{d}x}g(x)+p(x)\frac{\mathrm{d}g(x)}{\mathrm{d}x}.
\label{align:ch20-乘积导数}
\end{align}
对\eqref{align:ch20-乘积导数}的两边关于 $x$ 进行积分并整理，得出分部积分公式：
\begin{align}
\int_{-\infty}^{\infty}\frac{\mathrm{d}p(x)}{\mathrm{d}x}g(x)\,\mathrm{d}x=-\int_{-\infty}^{\infty}\frac{\mathrm{d}g(x)}{\mathrm{d}x}p(x)\,\mathrm{d}x.
\label{align:ch20-分部积分}
\end{align}
假设 $p(\infty)=p(-\infty)=0$。将此结果与定义 $\bm{s}_p=\nabla\ln p(\bm{x})$ 一起应用到涉及 $\bm{s}(\bm{x},\bm{w})^{\mathsf T}\bm{s}_p$ 的项以完成证明。请注意，计算\eqref{align:ch20-散度}中的二阶导数时，需要对每个导数进行一次反向传播，因此总的计算成本会随数据空间维度 $D$ 呈二次方增长（Martens, Sutskever, and Swersky, 2012）。因此，该损失函数不能直接应用于高维空间，于是切片得分匹配（Song et al., 2019）等技术被开发出来以解决这种低效问题。
\end{exercise}

\begin{exercise}\label{exercise:ch20-17}
（★★）证明：忽略加性常数，得分函数损失［\eqref{align:ch20-等价损失}］与\eqref{align:ch20-平滑损失}是等价的。首先展开\eqref{align:ch20-平滑损失}中的平方项，然后利用\eqref{align:ch20-Parzen}证明\eqref{align:ch20-平滑损失}中的 $\bm{s}^{\mathsf T}\bm{s}$ 项与通过展开\eqref{align:ch20-等价损失}中的平方项所得到的相应项是相同的。另请注意，\eqref{align:ch20-平滑损失}中 $\|\nabla_{\bm{z}}\ln q_{\sigma}\|^2$ 这一项与 $\bm{w}$ 无关，同样，\eqref{align:ch20-等价损失}中的相应项也与 $\bm{w}$ 无关，因此这些项可以视为损失函数中的加性常数，它们在训练中不起作用。最后考虑\eqref{align:ch20-平滑损失}中的交叉项。通过使用\eqref{align:ch20-Parzen}替换 $q_{\sigma}(\bm{z})$，证明这等于\eqref{align:ch20-等价损失}中相应的交叉项，从而证明这两个损失函数在忽略加性常数的情况下是相等的。
\end{exercise}

\begin{exercise}\label{exercise:ch20-18}
（★）考虑一个概率分布，它由两个互斥的分布（换言之，当其中一个非零时，另一个必定为零）混合而成，形式为
\begin{align}
p(\bm{x})=\lambda p_a(\bm{x})+(1-\lambda)p_b(\bm{x}).
\label{align:ch20-混合分布}
\end{align}
证明当针对任何给定点 $\bm{x}$ 计算由\eqref{align:ch20-得分}定义的得分函数时，混合系数 $\lambda$ 不会出现。由此可知，由\eqref{align:ch14-langevin更新}定义的 Langevin 动力学方法不会以正确的比例从这两个分布中采样。如正文所述，我们可以通过添加来自宽泛分布的噪声来解决这个问题。
\end{exercise}

\begin{exercise}\label{exercise:ch20-19}
（★★）在离散步骤，扩散模型中的正向加噪过程由\eqref{align:ch20-递推}定义。这里我们取连续时间极限，并将其转换为随机微分方程（SDE）。为此，首先引入一个连续变化的方差函数 $\beta(t)$，使得 $\beta_t=\beta(t)\Delta t$。通过对\eqref{align:ch20-递推}右边第一项中的平方根进行 Taylor 展开，证明无穷小的更新可以写成如下形式：
\begin{align}
\mathrm{d}\bm{z}=-\frac{1}{2}\beta(t)\,\bm{z}\,\mathrm{d}t+\sqrt{\beta(t)}\,\mathrm{d}\bm{w},
\label{align:ch20-无穷小}
\end{align}
这是一般随机微分方程［\eqref{align:ch20-SDE}］的特殊情况。
\end{exercise}

\begin{exercise}\label{exercise:ch20-20}
（★）用\eqref{align:ch20-条件得分}替换 $\nabla_{\bm{x}}\ln p(\bm{c}\mid\bm{x})$，证明\eqref{align:ch20-分类器引导}中的得分函数可以写成\eqref{align:ch20-无分类器引导}的形式。
\end{exercise}

\end{document}
